% ---------------------------------------------------------------------------
% Appendix: qualitative results across identities
% Figures: qual_identity_overview_id3.pdf (man), _id2.pdf (woman),
%          _id4.pdf (woman). One combined figure per identity: all variants
% side by side (like the main overview figure), BOTH views, mouth/eye insets,
% error map. UNIFORM 170-epoch budget (the fused TC-off arm has no 300-epoch
% run; figures never mix budgets). Per-clip dB values are view 0, frame 5.
% ---------------------------------------------------------------------------

\section{Qualitative Results Across Identities}
\label{sec:appendix_identities}

The qualitative figures in the main paper show a single validation subject.
To verify that the failure mode of joint view--temporal compression and its
recovery with latent capacity are not specific to one face, we repeat the
full comparison on three additional validation identities spanning gender and
skin tone (Figs.~\ref{fig:qual_identities_a}--\ref{fig:qual_identities_c}).
Each figure shows all model variants side by side at the 170-epoch budget of
the main comparison: the per-view and fused references without temporal
compression, the jointly compressed model at 16 channels, its widened 32- and
64-channel versions, and the best combination; rows show both camera views,
mouth and eye insets, and an amplified error map.

The pattern is consistent across identities. The per-view and fused
TC-off references reconstruct every subject cleanly (30--33\,dB on the shown
frames). Under joint view--temporal compression, degradation concentrates on
the articulated regions: teeth are smeared into a single band, mouth corners
are displaced, and eye detail is lost, while static regions (forehead,
cheeks, clothing) remain comparatively clean. Widening the latent and the
best combination recover much of the gap, in line with the quantitative
tables. The severity of the failure scales with articulation: for the
subject with a wide-open laugh (Fig.~\ref{fig:qual_identities_a}), the
16-channel jointly compressed reconstruction drops to 22.7\,dB with the mouth
interior essentially erased, whereas subjects with moderate expressions
degrade more mildly (26--27\,dB). The error maps confirm that residual error
concentrates on hair boundaries and facial articulation for every identity,
i.e.\ the compression bottleneck spends its budget on exactly the regions
that matter most for talking-head applications.

\begin{figure*}[t]
  \centering
  \includegraphics[width=\linewidth]{figures/qual_identity_overview_id3.pdf}
  \caption{All model variants for a male subject with a wide-open laugh
  (170 epochs). Joint view--temporal compression (fourth column) erases the
  mouth interior (22.7\,dB); latent capacity and the best combination
  progressively restore it.}
  \label{fig:qual_identities_a}
\end{figure*}

\begin{figure*}[t]
  \centering
  \includegraphics[width=\linewidth]{figures/qual_identity_overview_id2.pdf}
  \caption{All model variants for a second female subject (170 epochs).
  Teeth and eye detail degrade under joint compression and recover with
  capacity; skin tone is preserved by all variants.}
  \label{fig:qual_identities_b}
\end{figure*}

\begin{figure*}[t]
  \centering
  \includegraphics[width=\linewidth]{figures/qual_identity_overview_id4.pdf}
  \caption{All model variants for a third female subject (170 epochs). The
  same failure/recovery pattern holds; the teeth band and smile-narrowed
  eyes are the regions most affected by joint compression.}
  \label{fig:qual_identities_c}
\end{figure*}
